% GENERATED FILE. Edit canonical lessons, then run npm run generate:books.
% canonical-source-hash: fnv1a64:fb30265c2a09d22b
% canonical-lessons: TE-C178-kavalani, TE-C178-porapatuna, TE-C178-poni, TE-C178-saratu, TE-C178-saku

\chapter{On purpose, By mistake, Never mind, Condition, Excuse}
\label{ch:te-on-purpose-by-mistake-never-mind-condition-excuse}
\hlchaptermodality{178}

\begin{chapteropening}
\textbf{By the end of this chapter:} \emph{I can say on purpose, by mistake, never mind, a condition and an excuse.}

Complete the last lesson of chapter 178: I can say on purpose, by mistake, never mind, a condition and an excuse.
\end{chapteropening}

\section[\texorpdfstring{\te{కావాలని} (kāvālani)}{కావాలని (kāvālani)}]{\te{కావాలని} (kāvālani) — on purpose, deliberately}

\label{lesson:TE-C178-kavalani}



\begin{quote}
Before the new one: say the Telugu for actually, then the Telugu for anyway.
\end{quote}

\subsection*{You'll want to know: \te{కావాలని}}
\textbf{\te{కావాలని}} — \emph{kāvālani} — \textquotedblleft{}on purpose, deliberately\textquotedblright{}.

\begin{grammarlens}[title={What you've built}]
One more word for this chapter.
\end{grammarlens}

\subsection*{Your turn}
\begin{itemize}
\raggedright
  \item \emph{Say it:} \emph{kāvālani}
  \item \emph{Say it:} \emph{kāvālani}, once more
  \item \emph{Say it:} say \emph{kāvālani}
  \item \emph{Recall:} say the Telugu for actually, then the Telugu for anyway, then say \emph{kāvālani} again
\end{itemize}

\begin{tcolorbox}[breakable,colback=teal!4,colframe=teal!35!black,title={Before you move on}]
What does \te{కావాలని} mean? (\textquotedblleft{}On purpose, deliberately\textquotedblright{}.) Say it once more.
\end{tcolorbox}
\section[\texorpdfstring{\te{పొరపాటున} (porapāṭuna)}{పొరపాటున (porapāṭuna)}]{\te{పొరపాటున} (porapāṭuna) — by mistake, accidentally}

\label{lesson:TE-C178-porapatuna}



\begin{quote}
Before the new one: say the Telugu for anyway, then the Telugu for on purpose.
\end{quote}

\subsection*{You'll want to know: \te{పొరపాటున}}
\textbf{\te{పొరపాటున}} — \emph{porapāṭuna} — \textquotedblleft{}by mistake, accidentally\textquotedblright{}.

\begin{grammarlens}[title={What you've built}]
One more word for this chapter.
\end{grammarlens}

\subsection*{Your turn}
\begin{itemize}
\raggedright
  \item \emph{Say it:} \emph{porapāṭuna}
  \item \emph{Say it:} \emph{porapāṭuna}, once more
  \item \emph{Say it:} say \emph{porapāṭuna}
  \item \emph{Recall:} say the Telugu for anyway, then the Telugu for on purpose, then say \emph{porapāṭuna} again
\end{itemize}

\begin{tcolorbox}[breakable,colback=teal!4,colframe=teal!35!black,title={Before you move on}]
What does \te{పొరపాటున} mean? (\textquotedblleft{}By mistake, accidentally\textquotedblright{}.) Say it once more.
\end{tcolorbox}
\section[\texorpdfstring{\te{పోనీ} (pōnī)}{పోనీ (pōnī)}]{\te{పోనీ} (pōnī) — never mind; let it be}

\label{lesson:TE-C178-poni}



\begin{quote}
Before the new one: say the Telugu for on purpose, then the Telugu for by mistake.
\end{quote}

\subsection*{You'll want to know: \te{పోనీ}}
\textbf{\te{పోనీ}} — \emph{pōnī} — \textquotedblleft{}never mind; let it be\textquotedblright{}.

\begin{grammarlens}[title={What you've built}]
One more word for this chapter.
\end{grammarlens}

\subsection*{Your turn}
\begin{itemize}
\raggedright
  \item \emph{Say it:} \emph{pōnī}
  \item \emph{Say it:} \emph{pōnī}, once more
  \item \emph{Say it:} say \emph{pōnī}
  \item \emph{Recall:} say the Telugu for on purpose, then the Telugu for by mistake, then say \emph{pōnī} again
\end{itemize}

\begin{tcolorbox}[breakable,colback=teal!4,colframe=teal!35!black,title={Before you move on}]
What does \te{పోనీ} mean? (\textquotedblleft{}Never mind; let it be\textquotedblright{}.) Say it once more.
\end{tcolorbox}
\section[\texorpdfstring{\te{షరతు} (ṣaratu)}{షరతు (ṣaratu)}]{\te{షరతు} (ṣaratu) — a condition}

\label{lesson:TE-C178-saratu}



\begin{quote}
Before the new one: say the Telugu for by mistake, then the Telugu for never mind.
\end{quote}

\subsection*{You'll want to know: \te{షరతు}}
\textbf{\te{షరతు}} — \emph{ṣaratu} — \textquotedblleft{}a condition\textquotedblright{}.

\begin{grammarlens}[title={What you've built}]
One more word for this chapter.
\end{grammarlens}

\subsection*{Your turn}
\begin{itemize}
\raggedright
  \item \emph{Say it:} \emph{ṣaratu}
  \item \emph{Say it:} \emph{ṣaratu}, once more
  \item \emph{Say it:} say \emph{ṣaratu}
  \item \emph{Recall:} say the Telugu for by mistake, then the Telugu for never mind, then say \emph{ṣaratu} again
\end{itemize}

\begin{tcolorbox}[breakable,colback=teal!4,colframe=teal!35!black,title={Before you move on}]
What does \te{షరతు} mean? (\textquotedblleft{}A condition\textquotedblright{}.) Say it once more.
\end{tcolorbox}
\section[\texorpdfstring{\te{సాకు} (sāku)}{సాకు (sāku)}]{\te{సాకు} (sāku) — an excuse, a pretext}

\label{lesson:TE-C178-saku}



\begin{quote}
Before the new one: say the Telugu for never mind, then the Telugu for a condition.
\end{quote}

\subsection*{You'll want to know: \te{సాకు}}
\textbf{\te{సాకు}} — \emph{sāku} — \textquotedblleft{}an excuse, a pretext\textquotedblright{}.

\begin{grammarlens}[title={What you've built}]
One more word for this chapter.
\end{grammarlens}

\subsection*{Your turn}
\begin{itemize}
\raggedright
  \item \emph{Say it:} \emph{sāku}
  \item \emph{Say it:} \emph{sāku}, once more
  \item \emph{Say it:} say \emph{sāku}
  \item \emph{Recall:} say the Telugu for never mind, then the Telugu for a condition, then say \emph{sāku} again
\end{itemize}

\begin{tcolorbox}[breakable,colback=teal!4,colframe=teal!35!black,title={Before you move on}]
What does \te{సాకు} mean? (\textquotedblleft{}An excuse, a pretext\textquotedblright{}.) Say it once more.
\end{tcolorbox}
